\chapter{DISEÑO CONCEPTUAL Y DESCRIPCIÓN DEL ROBOT DIGITAL}
\label{ch:robot_digital}

Este capítulo describe la implementación del gemelo digital del robot móvil diferencial en \gls{ROS2} Foxy mediante URDF/\texttt{xacro} y Gazebo. Sobre la base del propósito de validación expuesto en la Sec.~\ref{sec:gemelo_digital}, se presentan su estructura, su backend de control y la correspondencia de parámetros con el prototipo físico.

En el código del proyecto, el archivo \texttt{robot.urdf.xacro} actúa como modelo principal del robot e integra los módulos de chasis, ruedas, sensor láser y control. El radio de rueda y la separación entre ruedas de la Tabla~\ref{tab:parametros_modelo} se emplean como referencia geométrica común en el robot físico, el URDF/\texttt{xacro} y \texttt{diff\_drive\_controller}. En cambio, la resolución de encoders y las frecuencias \texttt{loop\_rate}/\texttt{PID\_RATE} pertenecen a la ruta de hardware real: el backend \texttt{GazeboSystem} entrega directamente los estados de las articulaciones y no utiliza la conversión por ticks ni el planificador PID del firmware.

\section{Arquitectura del modelo digital}
\label{sec:atribucion}

El modelo digital se organiza mediante archivos \texttt{xacro} modulares. El archivo principal \texttt{robot.urdf.xacro} incluye la estructura mecánica del robot, definida en \texttt{robot\_core.xacro}, y selecciona el modelo de sensor según el entorno de ejecución: \texttt{lidar.xacro} para hardware real y \texttt{lidar\_sim.xacro} para simulación.

La contribución principal en esta etapa corresponde al diseño del modelo URDF/\texttt{xacro}, la parametrización geométrica del robot, la definición de marcos de referencia y la integración con los sistemas de control, SLAM y navegación. La atribución completa de paquetes propios, adaptados y externos se documenta en el \ref{anx:gemelo_digital}.

\section{Estructura URDF/xacro y marcos de referencia}
\label{sec:tf-convenciones}

El modelo se organiza como un árbol de eslabones y articulaciones cuyo marco principal es \texttt{base\_link}. En simulación, \texttt{slam\_toolbox} y AMCL utilizan \texttt{base\_footprint}, mientras que en hardware real se utiliza \texttt{base\_link}. La relación exacta entre ambos marcos y la verificación estructural del árbol se documentan en el \ref{anx:gemelo_digital}.

Se adoptan las convenciones REP-103 y REP-105: eje \texttt{x} hacia adelante, eje \texttt{y} hacia la izquierda y eje \texttt{z} hacia arriba \citep{rep103,rep105}. La cadena de transformaciones principal es:

\begin{center}
\texttt{map} $\rightarrow$ \texttt{odom} $\rightarrow$ \texttt{base\_link} $\rightarrow$ \texttt{laser\_frame},
\end{center}

\noindent con la rama fija y coincidente \texttt{base\_link}$\rightarrow$\texttt{base\_footprint}.

\begin{table}[H]
\centering
\caption{Marcos y tópicos principales del modelo digital.}
\label{tab:topics_frames}
\small
\begin{tabular}{p{3.5cm} p{4.0cm} p{5.0cm}}
\toprule
\textbf{Componente} & \textbf{Marco} & \textbf{Tópico asociado} \\
\midrule
Base diferencial & \texttt{odom} $\rightarrow$ \texttt{base\_link} & \texttt{/odom} \\
Sensor láser & \texttt{laser\_frame} & \texttt{/scan} \\
Transformación estática & \texttt{base\_link} $\rightarrow$ \texttt{laser\_frame} & \texttt{/tf\_static} \\
\bottomrule
\end{tabular}
\end{table}

La validación estructural del modelo, incluyendo el conteo de eslabones, articulaciones y árbol TF completo, se traslada al \ref{anx:gemelo_digital}.

\section{Selección del backend de control}
\label{sec:backend}

El modelo permite seleccionar el backend mediante los argumentos \texttt{use\_ros2\_control} y \texttt{sim\_mode}. La ejecución utilizada mantiene \texttt{use\_ros2\_control=true} en ambos entornos: el robot real carga \texttt{diffdrive\_arduino/DiffDriveArduino} y Gazebo carga \texttt{gazebo\_ros2\_control/GazeboSystem}. La trazabilidad del backend alternativo no ejecutado se conserva en el \ref{anx:gemelo_digital}.

\begin{table}[H]
\centering
\caption{Backend de control según entorno de ejecución.}
\label{tab:xacro_args_backend}
\small
\begin{tabular}{p{3.0cm} p{3.0cm} p{6.0cm}}
\toprule
\textbf{Entorno} & \textbf{Backend} & \textbf{Descripción} \\
\midrule
Hardware real & \texttt{ros2\_control} & Plugin \texttt{DiffDriveArduino} y comunicación serial con el \gls{Arduino}. \\
Simulación & \texttt{ros2\_control} & Plugin \texttt{gazebo\_ros2\_control/GazeboSystem} cargado por \texttt{libgazebo\_ros2\_control.so}. \\
\bottomrule
\end{tabular}
\end{table}

\section{Sensor láser simulado}

En simulación, el sensor se modela mediante \texttt{lidar\_sim.xacro}, publicando datos en \texttt{/scan} con el marco \texttt{laser\_frame}. El modelo reproduce el comportamiento general del RPLidar A1 físico, aunque con diferencias propias de la simulación, como ausencia de ruido real y un rango mínimo más permisivo. Estas diferencias se consideran parte de la brecha simulación--realidad.

\section{Criterios de aceptación}

Antes de usar el gemelo digital para pruebas de SLAM y navegación, se verifican tres aspectos mínimos: coherencia de la cadena TF, publicación de tópicos principales y respuesta cinemática ante comandos de velocidad.

Los criterios considerados fueron la continuidad de la cadena TF, la
publicación de \texttt{/scan} y \texttt{/odom}, la respuesta cinemática y la
paridad geométrica con el prototipo. La tabla de verificación y la evidencia
estructural se trasladan al \ref{anx:gemelo_digital},
Tabla~\ref{tab:criterios_aceptacion_resumen}.

\section*{Conclusión del capítulo}

El gemelo digital permite validar de forma previa la geometría, la cinemática, los marcos de referencia, los tópicos y la configuración de control del robot diferencial. Al mantener parámetros compartidos entre Gazebo y el hardware real, se facilita la transferencia hacia el prototipo físico y se establece una base reproducible para las pruebas de SLAM y navegación.